% ============================================================================
% Fichier  : partie1/P01_C05_grandes_attaques.tex
% Titre    : Histoire des Grandes Attaques
% Auteur   : MINKA MI NGUIDJOÏ Thierry Emmanuel
% Version  : 1.0 -- novembre 2026
% ============================================================================
% Base : brouillon, introduction (historique). Sources : Spafford (1989) pour
% le ver Morris ; Falliere, Murchu et Chien (2011) pour Stuxnet ; rapport du
% National Audit Office (2017) pour WannaCry ; Greenberg (2019) pour
% NotPetya ; alerte CISA AA20-352A et déclaration 8-K de SolarWinds (2020).
% Aucun chiffre n'est donné sans source ; les estimations sont signalées.
% ============================================================================

\chapter{Histoire des Grandes Attaques}
\label{chap:attaques}
\chaptermark{Grandes attaques}

\epigraph{%
    « Les attaques changent d'outils plus vite que de méthode. »%
}{Principe du cours}

\niveauChapitre{Fondamental}{Chapitres~\ref{chap:proprietes}
à~\ref{chap:crypto-asym}.}

\begin{boxobjectifs}
\begin{itemize}
    \item Lire une attaque historique avec une grille constante.
    \item Distinguer source primaire, source secondaire et estimation.
    \item Décrire cinq attaques majeures et leurs mécanismes.
    \item Dégager les leçons communes et les relier aux mesures du cours.
    \item Transposer ces leçons au contexte camerounais.
\end{itemize}
\end{boxobjectifs}

\bigskip

% ============================================================================
\section*{Question directrice : l'Opération GAROUA était-elle prévisible~?}
% ============================================================================

Chacun des maillons de l'Opération GAROUA a un précédent célèbre~: un
correctif publié et non installé, comme pour WannaCry~; un prestataire de
confiance utilisé comme porte d'entrée, comme pour NotPetya et
SolarWinds~; un rançongiciel qui détruit plus qu'il ne rançonne. Étudier
l'histoire, c'est apprendre à reconnaître ce qui revient.

\medskip
\textbf{Plan du chapitre.}
\begin{enumerate}
    \item Une grille de lecture (section~\ref{sec:c05-grille}).
    \item Le ver Morris, 1988 (section~\ref{sec:c05-morris}).
    \item Stuxnet, 2010 (section~\ref{sec:c05-stuxnet}).
    \item WannaCry et NotPetya, 2017 (section~\ref{sec:c05-2017}).
    \item SolarWinds, 2020 (section~\ref{sec:c05-solarwinds}).
    \item Les leçons communes (section~\ref{sec:c05-lecons}).
\end{enumerate}

\separateur

% ============================================================================
\section{Une Grille de Lecture}
\label{sec:c05-grille}
% ============================================================================

Une attaque se lit avec les mêmes questions, quelle qu'elle soit~:

\begin{tableaucours}{L{3.4cm}L{8.8cm}}{%
    \tableauHeader{Question} &
    \tableauHeader{Ce qu'elle cherche}
}
Les faits & Quoi, quand, qui a été touché, d'après quelles sources \\
Le vecteur & Par où l'attaque est-elle entrée~? \\
Le mécanisme & Comment s'est-elle propagée et qu'a-t-elle fait~? \\
La mesure manquante & Quelle mesure l'aurait arrêtée ou limitée~? \\
Le signal ignoré & Qu'aurait-on pu voir avant~? \\
La leçon & Qu'a-t-elle changé dans les pratiques ou les textes~? \\
\end{tableaucours}
\vspace{-0.8em}
\begin{center}{\small\itshape Tableau~5.1. La grille de lecture d'une attaque}\end{center}

\begin{boxalert}
\textbf{Violation typique.} Citer le coût d'une attaque lu dans un article
de presse comme un fait établi.

\smallskip
\textbf{Règle~:} on distingue la \emph{source primaire} (le rapport de
l'organisme touché, d'une autorité, d'un tribunal), la \emph{source
secondaire} (presse, ouvrage) et l'\emph{estimation}, que l'on présente
comme telle. Les montants des grandes attaques sont presque toujours des
estimations.
\end{boxalert}

% ============================================================================
\section{Le Ver Morris, 1988}
\label{sec:c05-morris}
% ============================================================================

Le 2~novembre~1988, un programme écrit par un étudiant de Cornell, Robert
T.~Morris, se propage sur l'Internet de l'époque. E.~Spafford, qui l'a
analysé, décrit trois vecteurs~: un mode de débogage laissé actif dans le
serveur de messagerie \texttt{sendmail}, un débordement de tampon dans le
service \texttt{fingerd}, et des mots de passe faibles ou des relations de
confiance entre machines~\cite{spafford1989}. Une erreur de conception le
fait se réinstaller sur des machines déjà infectées, qui s'effondrent sous
la charge. Le nombre de machines touchées n'a jamais été mesuré~; les
estimations de l'époque parlent de quelques milliers.

\textbf{Ce qu'il a changé.} La création, quelques jours plus tard, du
premier centre de réponse aux incidents (le CERT/CC à Carnegie Mellon), et
la première condamnation aux États-Unis au titre de la loi fédérale sur la
fraude informatique.

% ============================================================================
\section{Stuxnet, 2010}
\label{sec:c05-stuxnet}
% ============================================================================

Découvert en 2010, Stuxnet vise des automates industriels Siemens qui
pilotent des centrifugeuses. N.~Falliere, L.~O~Murchu et E.~Chien en ont
publié l'analyse détaillée~\cite{falliere2011}~: propagation par clé USB
pour franchir l'isolement des réseaux industriels, exploitation de
plusieurs vulnérabilités inconnues jusque-là, pilotes signés avec des
certificats volés, et modification discrète de la vitesse des moteurs
pendant que les opérateurs voyaient des valeurs normales.

\begin{boxinfo}
\textbf{Ce que Stuxnet enseigne.} Un réseau isolé n'est pas un réseau
protégé~: une clé USB suffit. Une signature numérique vaut ce que vaut la
protection de la clé privée qui l'a produite (Chapitre~\ref{chap:crypto-asym}).
Et une attaque peut viser l'intégrité d'un processus physique, pas
seulement des données.
\end{boxinfo}

% ============================================================================
\section{WannaCry et NotPetya, 2017}
\label{sec:c05-2017}
% ============================================================================

\subsection{WannaCry}

Le 12~mai~2017, le rançongiciel WannaCry se propage seul, comme un ver, en
exploitant une vulnérabilité du protocole de partage de fichiers SMBv1 de
Windows, pour laquelle Microsoft avait publié un correctif (MS17-010) en
mars. Au Royaume-Uni, le National Audit Office a établi que des
établissements du service public de santé ont été touchés, des rendez-vous
annulés, et que la plupart des machines infectées n'avaient pas reçu le
correctif~\cite{nao2017wannacry}. La propagation s'arrête quand un
chercheur enregistre un nom de domaine que le programme interrogeait avant
d'agir.

\subsection{NotPetya}

Le 27~juin~2017, un programme qui se présente comme un rançongiciel se
diffuse à partir de la mise à jour piégée d'un logiciel de comptabilité
ukrainien. Il chiffre de façon irréversible~: c'est une arme de
destruction déguisée. Il se propage dans les réseaux des entreprises en
réutilisant des identifiants récupérés en mémoire et la vulnérabilité
exploitée par WannaCry~; de grandes entreprises internationales voient
leur informatique arrêtée~\cite{greenberg2019}.

\begin{boxscenario}{Opération GAROUA. Les échos de 2017}
\textbf{Comme WannaCry}~: un correctif critique, publié huit mois plus
tôt, n'était pas installé sur le serveur de fichiers. \textbf{Comme
NotPetya}~: l'attaque est entrée par un acteur de confiance (le
prestataire) et s'est propagée grâce à des identifiants d'administrateur
réutilisés sur plusieurs machines.

\medskip
\textbf{Analyse} (confiance élevée). Les mesures étaient connues depuis
2017~: correctifs, segmentation, comptes d'administration distincts par
machine, sauvegardes hors ligne.
\end{boxscenario}

% ============================================================================
\section{SolarWinds, 2020}
\label{sec:c05-solarwinds}
% ============================================================================

En décembre~2020, on découvre qu'une mise à jour du logiciel de
supervision Orion, de l'éditeur SolarWinds, contenait une porte dérobée
insérée dans la chaîne de production du logiciel. L'agence américaine de
cybersécurité (CISA) a publié une alerte décrivant l'attaque et ses
indicateurs~\cite{cisa2020solarwinds}~; l'éditeur a déclaré à l'autorité
boursière qu'environ 18\,000 clients avaient pu installer la version
compromise~\cite{solarwinds2020}, l'attaquant n'en exploitant ensuite
qu'une petite partie.

\begin{boxinfo}
\textbf{Ce que SolarWinds enseigne.} La confiance dans un fournisseur est
un risque. Une mise à jour signée par l'éditeur peut être malveillante si
la chaîne de production de l'éditeur est compromise. D'où l'intérêt de
surveiller le comportement d'un logiciel, même de confiance, et de limiter
ses droits (Chapitre~\ref{chap:architecture}).
\end{boxinfo}

% ============================================================================
\section{Les Leçons Communes}
\label{sec:c05-lecons}
% ============================================================================

\begin{tableaucours}{L{2.6cm}L{3.2cm}L{3.2cm}L{3.2cm}}{%
    \tableauHeader{Attaque} &
    \tableauHeader{Vecteur} &
    \tableauHeader{Mesure manquante} &
    \tableauHeader{Chapitre du cours}
}
Morris & Services mal configurés, mots de passe faibles & Durcissement,
mots de passe & \ref{chap:systemes}, \ref{chap:iam} \\
Stuxnet & Clé USB, vulnérabilités inconnues & Contrôle des supports,
détection & \ref{chap:actifs}, \ref{chap:operations} \\
WannaCry & Faille SMBv1 corrigée depuis deux mois & Correctifs,
segmentation & \ref{chap:systemes}, \ref{chap:architecture} \\
NotPetya & Mise à jour piégée d'un fournisseur & Segmentation, comptes
distincts, sauvegardes & \ref{chap:architecture}, \ref{chap:continuite} \\
SolarWinds & Chaîne de production d'un éditeur & Moindre privilège,
surveillance & \ref{chap:architecture}, \ref{chap:operations} \\
\end{tableaucours}
\vspace{-0.8em}
\begin{center}{\small\itshape Tableau~5.2. Cinq attaques, une grille}\end{center}

Quatre leçons reviennent~: \textbf{(1)} la plupart des attaques exploitent
des faiblesses connues~; \textbf{(2)} la confiance accordée à un tiers est
un vecteur~; \textbf{(3)} un réseau plat transforme une intrusion locale en
catastrophe générale~; \textbf{(4)} la capacité à reprendre l'activité
(sauvegardes, plan de continuité) décide de l'issue.

\begin{boxinfo}
\textbf{Et en Afrique~?} Les attaques survenues en Afrique centrale sont
rarement documentées par des sources primaires publiques. Le cours ne
cite que des faits dont la source est accessible~; l'exercice~5.10
invite à en rechercher, avec la même exigence.
\end{boxinfo}

\begin{boxmetier}{Auditeur de sécurité}
\textbf{Se servir de l'histoire.} L'auditeur transforme chaque leçon en
point de contrôle~: les correctifs critiques sont-ils installés sous
trente jours~? Les comptes d'administration sont-ils distincts par
machine~? Les sauvegardes sont-elles hors ligne et testées~? Les
fournisseurs ont-ils des droits limités~?
\end{boxmetier}

\separateur

\begin{boxresume}
\textbf{Ce qu'il faut retenir de ce chapitre~:}
\begin{itemize}
    \item \textbf{Une attaque se lit avec une grille}~: faits, vecteur,
          mécanisme, mesure manquante, signal ignoré, leçon.
    \item \textbf{Les sources se distinguent}~; les montants sont presque
          toujours des estimations.
    \item \textbf{Morris, Stuxnet, WannaCry, NotPetya, SolarWinds}
          illustrent les mots de passe faibles, l'isolement illusoire, le
          correctif manquant, le fournisseur piégé.
    \item \textbf{Quatre leçons}~: faiblesses connues, confiance dans les
          tiers, réseau plat, capacité de reprise.
\end{itemize}
\end{boxresume}

\begin{boxinfo}
\textbf{Points de contrôle de l'élève.} Je sais~: appliquer la grille à une
attaque~; distinguer source primaire et estimation~; décrire le mécanisme
des cinq attaques~; associer chaque attaque à une mesure du cours.
\end{boxinfo}

\bigskip

\begin{boxexo}[title={\ding{45}\quad Exercice~5.1. La grille \niveaubase}]
Appliquez la grille du Tableau~5.1 à l'Opération GAROUA.
\end{boxexo}

\begin{boxexo}[title={\ding{45}\quad Exercice~5.2. Quelle attaque~? \niveaubase}]
Associez chaque phrase à une attaque~: (a) une clé USB franchit un réseau
isolé~; (b) un correctif publié en mars n'est pas installé en mai~; (c)
une mise à jour de logiciel de comptabilité détruit des réseaux entiers~;
(d) un programme se réinstalle sans fin sur les mêmes machines.
\end{boxexo}

\begin{boxexo}[title={\ding{45}\quad Exercice~5.3. Les sources \niveaubase}]
Classez en source primaire, secondaire ou estimation~: le rapport du
National Audit Office~; un article qui chiffre le coût mondial de
NotPetya~; la déclaration 8-K de SolarWinds~; un billet de blogue sur
Stuxnet.
\end{boxexo}

\begin{boxexo}[title={\ding{45}\quad Exercice~5.4. Le correctif \niveaumoyen}]
Entre la publication d'un correctif et son installation, une entreprise
est exposée. Proposez une règle de délai par niveau de gravité, et la
mesure compensatoire quand un correctif ne peut pas être installé.
\end{boxexo}

\begin{boxexo}[title={\ding{45}\quad Exercice~5.5. L'isolement \niveaumoyen}]
La SODECO, entreprise fictive de distribution d'électricité, pilote ses
postes par un réseau industriel « déconnecté d'Internet ». Quelles leçons
de Stuxnet s'y appliquent~?
\end{boxexo}

\begin{boxexo}[title={\ding{45}\quad Exercice~5.6. Le fournisseur \niveaumoyen}]
LiZbiZ utilise un logiciel de paie mis à jour automatiquement par son
éditeur. Quelles mesures limitent le risque d'une mise à jour piégée~?
\end{boxexo}

\begin{boxexo}[title={\ding{45}\quad Exercice~5.7. Réseau plat \niveaumoyen}]
Expliquez pourquoi WannaCry et NotPetya ont été dévastateurs dans les
réseaux plats, et décrivez une segmentation minimale pour LiZbiZ.
\end{boxexo}

\begin{boxexo}[title={\ding{45}\quad Exercice~5.8. Rançongiciel ou destructeur \niveaumoyen}]
NotPetya ressemblait à un rançongiciel mais ne permettait pas de
récupérer les données. Qu'est-ce que cela change pour la décision de payer
une rançon~?
\end{boxexo}

\begin{boxexo}[title={\ding{45}\quad Exercice~5.9. Des points de contrôle \niveauexpert}]
À partir du Tableau~5.2, rédigez dix points de contrôle pour un audit de
sécurité de la BANQUE SANAGA, chacun avec la preuve attendue.
\end{boxexo}

\begin{boxexo}[title={\ding{45}\quad Exercice~5.10. Une affaire africaine \niveauexpert}]
Recherchez une attaque survenue en Afrique et documentée par une source
primaire publique (rapport d'une autorité, décision de justice,
communiqué de l'organisation touchée). Appliquez la grille, en signalant
le statut de chaque source.
\end{boxexo}

\bigskip

\begin{boxinfo}
\textbf{Pour aller plus loin.}
\begin{itemize}
    \item E.~Spafford, analyse du ver de 1988~\cite{spafford1989}.
    \item N.~Falliere, L.~O~Murchu, E.~Chien, \emph{W32.Stuxnet
          Dossier}~\cite{falliere2011}.
    \item National Audit Office, rapport sur WannaCry~\cite{nao2017wannacry}~;
          A.~Greenberg, \emph{Sandworm}~\cite{greenberg2019}.
\end{itemize}
\end{boxinfo}

\begin{boxinfo}
\textbf{Transition.} Ces attaques ont réussi autant par défaut
d'organisation que par défaut technique. La Partie~\ref{part:gouvernance}
commence par la gouvernance~: qui décide, qui est responsable, et selon
quelle politique.
\end{boxinfo}
